\subsection*{Who produced the labels} A single annotator applied the frozen
guideline, so the labels carry no inter-annotator agreement statistic. A blind
second pass over already-labelled items, shown in a different order with the
first-pass labels withheld, reproduced them exactly (27 of 27, $\kappa=1.00$). We
report the figure and rest little on it, since perfect self-agreement is as
consistent with an annotator recognising passages seen earlier as with a
guideline that admits one reading. The quantity we do rest on is the
within-annotator context ablation (\S\ref{sec:models}), which measures how far
the same annotator's decisions depend on material outside the anchor sentence.
Whether two trained lawyers would agree on these labels is a question for a
follow-up study with a practitioner panel.

The annotator is a language model, which places it in the same family as some of
the systems under evaluation. Two design choices contain the circularity. The
supervised baselines are trained from the labels but are architecturally
independent of the label source, and they carry the context-scaling result. The
zero-shot baseline is a small open model that had no part in producing the
labels, saw no guideline, and was scored by label-token likelihood, so it is
independent of them. A frontier model evaluated against its own annotations would
report agreement with itself, so we omit that comparison. The escalation results
in \S\ref{sec:long} rest on the same labels.

\subsection*{Benchmark and split sizes} Annotation produced \NBench{} usable
labels, \NDev{} for development and 104, 36, and 46 for the three held-out
conditions. At that size, differences of a few macro-$F_1$ points sit inside the
noise, and the temporal condition carries wide intervals. The two comparisons we
rest on are the per-stratum contrast, where the gap between the sparse model and
the dictionary on H1 is large, and the paired context comparison, which is
within-item and therefore better powered than the split sizes suggest. One fact a
reader should weigh is that the separation between supervised systems and the
dictionary emerged as the development split grew, having been absent at half this
annotation scale. That is the signature of an effect that needed data to surface,
and it is also a reason to rerun the benchmark at several thousand items.

The leakage control reaches \IssueOnlyF{} macro-$F_1$ against a majority baseline
of 0.392, with intervals that barely overlap. Every dictionary expression and
holding cue is stripped from a statement before it is kept, so the signal here is
subject matter rather than lexical leakage of the label. Which questions courts
decline to reach is correlated with what the questions are about, so part of what
a system appears to infer from the passage is available from the topic alone.
Future versions of this benchmark should either match strata on subject matter or
report the issue-only control alongside every headline number, as we do here.

\subsection*{What the dictionary measures} A court can leave an issue open in
ordinary prose. Freezing the thirteen expressions before annotation protects
against tuning them to a result, and the recall audit in Appendix
\ref{app:recall} bounds what they miss by reading a random sample of opinions
containing none of them. Prevalence figures should be read as the rate of
explicitly marked non-resolution. The item frame is scoped the same way. It
admits anchors with a recoverable issue statement, and yield varies across
expressions and, less strongly, across courts, so the benchmark samples anchors
with a recoverable proposition. Appendix \ref{app:issue} reports yield by
stratum so the direction of the loss is visible.

\subsection*{Corpus coverage} CourtListener is broad but uneven. Coverage,
metadata completeness, and text quality vary by court and era, and the same
opinion is sometimes stored more than once from different upstream sources. We
deduplicate and report the eligibility funnel, and a change in measured
prevalence across periods remains partly a change in what was digitized and how.
This is why period-stratified rates appear throughout rather than a single time
series presented as a trend in judicial behaviour.

A citing passage is treated as being about the origin's issue when its content
words overlap the issue statement. That admits false positives, where a later
opinion cites the origin for a neighbouring point sharing vocabulary, and false
negatives, where a later court restates the issue in different words. Annotation
removes the false positives, since off-issue passages are labelled and dropped,
and the false-negative rate is unmeasured.

\subsection*{Scope of the comparison} The comparison between unresolved and
decided issues matches on court, period, opinion length, and citation volume, and
is descriptive. Whether an issue was left open is a judicial choice correlated
with unobservables, including how hard and how contested the issue was, and those
plausibly affect how later courts describe it. We report the contrast as an
association.

Chains require the origin to carry a reporter citation locatable in the citing
text, and are restricted to citing opinions inside the population, so district
and state court treatments fall outside the sample and origins in the most recent
years have had little time to accumulate citations. Both restrictions bias the
chain sample toward older, better-reported, more-cited opinions.

The encoder baselines truncate at 512 tokens and the zero-shot model at 1600,
which caps the widest context condition, so the comparison across context widths
is partly a comparison of how much of the window each architecture uses. The
zero-shot model has 3.8B parameters, and its role here is to show that an
independent model responds to the context manipulation in the same direction. The
study covers English-language United States federal appellate opinions, and
extension to state courts, trial courts, agency adjudication, or other legal
systems remains future work.
